\documentclass[11pt]{amsart}

\usepackage[margin=1.15in]{geometry}
\usepackage{amsmath,amssymb,amsthm,mathtools}
\usepackage{microtype}
\usepackage{hyperref}
\usepackage[nameinlink,capitalize]{cleveref}

\hypersetup{
  colorlinks=true,
  linkcolor=blue,
  citecolor=blue,
  urlcolor=blue
}

\theoremstyle{plain}
\newtheorem{theorem}{Theorem}[section]
\newtheorem{proposition}[theorem]{Proposition}
\newtheorem{lemma}[theorem]{Lemma}
\newtheorem{corollary}[theorem]{Corollary}

\theoremstyle{definition}
\newtheorem{definition}[theorem]{Definition}

\theoremstyle{remark}
\newtheorem{remark}[theorem]{Remark}

\numberwithin{equation}{section}

\title[Uniformly tight ancient solutions]{Uniformly Tight Ancient Solutions of the Navier--Stokes Equations in Self-similar Variables}

\author{Guillem Duran Ballester}
\date{}

\subjclass[2020]{35Q30, 35B44, 76D05}
\keywords{Navier--Stokes equations, ancient solutions, self-similar variables, Type I singularities, Liouville theorems}

\begin{document}

\begin{abstract}
We study bounded mild ancient solutions of the three-dimensional
Navier--Stokes equations in backward self-similar variables.  We establish
local regularity, compactness, and stability of the Leray pressure under a
uniform \(L^\infty(\mathbb R_\tau;L^3(\mathbb R^3_y))\) bound.  We prove
that any locally smooth time-translate limit which is stationary and arises
from a uniformly \(L^3\)-tight ancient solution is identically zero.  The
proof invokes the stationary self-similar rigidity theorem of
Ne\v{c}as--R\r{u}\v{z}i\v{c}ka--\v{S}verak.
We also prove a small-data Liouville theorem for bounded mild ancient
solutions, using only the mild formulation and the decay of the
Ornstein--Uhlenbeck--Stokes semigroup.
\end{abstract}

\maketitle

\section{Introduction}

We consider the three-dimensional incompressible Navier--Stokes equations in
backward self-similar variables.  The velocity \(V\) and pressure \(P\) solve
\begin{equation}\label{eq:centered}
   \partial_\tau V+(V\cdot\nabla)V+\nabla P
   =
   \Delta V-\frac12(V+y\cdot\nabla V),
   \qquad \operatorname{div} V=0
\end{equation}
on \(\mathbb R^3\times\mathbb R\).  This equation is the autonomous form obtained from the
physical Navier--Stokes system by the change of variables
\[
   V(y,\tau)=\sqrt{-t}\,u(y\sqrt{-t},t),\qquad \tau=-\log(-t),
   \qquad t<0.
\]
The drift term fixes the self-similar scale.  The remaining scaling action is
translation in the logarithmic time variable \(\tau\).

Throughout the paper repeated spatial indices are summed, and \(R_i\) denotes
the \(i\)-th Riesz transform on \(\mathbb R^3\), with Fourier symbol
\(i\xi_i/|\xi|\).  Thus \(R_iR_j\) has symbol
\(-\xi_i\xi_j/|\xi|^2\), and
\[
   P=R_iR_j(F_{ij})
   \quad\Longleftrightarrow\quad
   -\Delta P=\partial_i\partial_jF_{ij}
\]
in distributions.  All differential equations are understood in the sense of
distributions unless a stronger sense is explicitly stated.  Pressures are
determined only modulo functions of time; a ``gauge'' always means the choice
of such a time-dependent additive function.  Constants denoted by \(C\) may
change from line to line, but their displayed dependencies are the only
dependencies used in the argument.

This paper does not assert a full ancient Liouville theorem for
\eqref{eq:centered}.  It isolates two conditional Liouville criteria which
are available under additional compactness or smallness hypotheses.

First, uniform tightness in \(L^3(\mathbb R^3)\) permits one to pass from
local smooth convergence to a global \(L^3\) stationary self-similar profile.
The theorem of Ne\v{c}as, R\r{u}\v{z}i\v{c}ka, and \v{S}verak
\cite{NRS1996} then implies that the stationary limit is zero.  Second, a
bounded mild ancient solution with sufficiently small
\(L^\infty(\mathbb R^3\times\mathbb R)\) norm is zero.  This follows from the
mild formulation of \eqref{eq:centered} and the decay of the linear
Ornstein--Uhlenbeck--Stokes semigroup.

\subsection{Logical scope}

The results below are separated according to their hypotheses.  The global
source-level condition
\(u\in L^\infty(0,T;L^3(\mathbb R^3))\) is ruled out by the
Escauriaza--Seregin--\v{S}verak theorem before any ancient-profile classification
is required.  The small-amplitude ancient regime is ruled out by the mild
formula and the damping of the centered Stokes semigroup.  The stationary
\(L^3\) regime is ruled out by the stationary self-similar theorem.  The
uniformly \(L^3\)-tight time-dependent regime is not excluded here unless an
additional argument produces a stationary limit.

The source-level facts needed below are stated here as analytic hypotheses on
the ancient solution under consideration.  Namely, the solution is assumed to
arise from a pressure-normalized CKN concentration sequence at a Type I
singular point, after admissible pressure gauges and scale selection, and to
carry the corresponding nonvanishing normalization.  These concentration,
gauge, compactness, and nontriviality requirements are precisely the
conditional input needed for the tight ancient dynamics argument.

\subsection{Mild formulation}

After applying the Leray projection \(\mathbb P\), equation
\eqref{eq:centered} becomes
\begin{equation}\label{eq:projected}
   \partial_\tau V=
   \Delta V-\frac12 y\cdot\nabla V-\frac12 V
   -\mathbb P\operatorname{div}(V\otimes V).
\end{equation}
The linear semigroup associated with the first three terms on the right-hand
side is
\begin{equation}\label{eq:OU-semigroup}
   (S(a)f)(y)
   =
   e^{-a/2}
   \int_{\mathbb R^3}
   \frac{\exp\!\left(-\frac{|z-e^{-a/2}y|^2}{4(1-e^{-a})}\right)}
        {[4\pi(1-e^{-a})]^{3/2}}
   f(z)\,dz,
   \qquad a>0.
\end{equation}
The factor \(e^{-a/2}\) is the damping coming from the term \(-V/2\).
Equivalently, \(S(a)=e^{aL}\), where
\[
   L=\Delta-\frac12y\cdot\nabla-\frac12 .
\]
The formula shows that \(S(a)\) maps bounded functions to bounded functions,
preserves distributional divergence-free vector fields, and satisfies the
semigroup law \(S(a)S(b)=S(a+b)\).  These facts follow from the change of
variables representation of the Ornstein--Uhlenbeck heat kernel and the
corresponding facts for the heat semigroup.

\begin{definition}[Bounded mild ancient solution]\label{def:mild}
A vector field \(V:\mathbb R^3\times\mathbb R\to\mathbb R^3\) is a bounded mild ancient solution of
\eqref{eq:centered} if
\[
   V\in L^\infty(\mathbb R^3\times\mathbb R),\qquad \operatorname{div} V=0
\]
in distributions, and if for every \(-\infty<\sigma<\tau<\infty\)
\begin{equation}\label{eq:mild}
   V(\tau)
   =
   S(\tau-\sigma)V(\sigma)
   -
   \int_\sigma^\tau
   S(\tau-s)\mathbb P\operatorname{div}(V\otimes V)(s)\,ds
\end{equation}
as an identity in the weak-* topology of \(L^\infty(\mathbb R^3)\).  A pressure
\(P\) is then recovered locally from
\[
   -\Delta P=\partial_i\partial_j(V_iV_j)
\]
modulo functions of time.
\end{definition}

The weak-* formulation means that, for every
\(\phi\in L^1(\mathbb R^3;\mathbb R^3)\), pairing \eqref{eq:mild} against
\(\phi\) gives a scalar identity.  The Duhamel term is interpreted by duality:
for compactly supported smooth \(\phi\) and \(s<\tau\),
\[
   \left\langle
      S(\tau-s)\mathbb P\operatorname{div}(V\otimes V)(s),\phi
   \right\rangle
   =
   \left\langle
      V\otimes V(s),
      (S(\tau-s)\mathbb P\operatorname{div})^*\phi
   \right\rangle .
\]
For fixed \(\sigma<\tau\), \cref{lem:stokes-estimates} gives
\[
   \bigl|
      \langle
      S(\tau-s)\mathbb P\operatorname{div}(V\otimes V)(s),\phi
      \rangle
   \bigr|
   \le
   C\|V\|_{L^\infty}^2
   e^{-(\tau-s)/2}(1-e^{-(\tau-s)})^{-1/2}\|\phi\|_{L^1},
\]
and the right-hand side is integrable for \(s\in(\sigma,\tau)\).  Hence the
time integral in \eqref{eq:mild} is a well-defined weak-* integral in
\(L^\infty(\mathbb R^3)\).  Once the integral has been represented by this
bounded function, equality against compactly supported smooth tests extends
by density of \(C_c^\infty(\mathbb R^3)\) in \(L^1(\mathbb R^3)\) to equality
against all \(L^1\) tests.  Since \(L^\infty(\mathbb R^3)\) is identified
with the dual of \(L^1(\mathbb R^3)\) modulo almost-everywhere equality, two
bounded representatives which have the same pairing against every
\(L^1\)-test agree almost everywhere.  Thus no pointwise pressure
representative is used in the definition of mild solution.

\begin{remark}
The mild condition fixes the whole-space pressure normalization and excludes
the local pressure ambiguities which are invisible to the projected equation.
For whole-space ancient limits obtained from Leray solutions, this is the
Leray pressure, modulo functions of time, rather than an arbitrary local
pressure with additional affine spatial components.  The invariance under
time-dependent pressure gauges in the source suitable-weak setting follows
directly because such gauges have zero spatial gradient and cancel from
pressure oscillations.
\end{remark}

\subsection{Main results}

The first result is a local regularity and pressure compactness statement.

\begin{theorem}[Regularity and pressure compactness]\label{thm:gauges}
Let \(V\) be a bounded mild ancient solution of \eqref{eq:centered}.  Then
\(V\) is smooth on \(\mathbb R^3\times\mathbb R\).  Moreover, if \(V_n\) is a sequence of
bounded mild ancient solutions satisfying
\[
   \sup_n\|V_n\|_{L^\infty(\mathbb R^3\times\mathbb R)}<\infty
\]
and, in addition,
\[
   \sup_n\sup_{\tau\in\mathbb R}\|V_n(\tau)\|_{L^3(\mathbb R^3)}<\infty,
\]
then the Leray pressures are stable under local limits.  More precisely, if
\(P_n=R_iR_j(V_{n,i}V_{n,j})\) are the Leray pressures and
\(V_n\to V\) locally uniformly, then
\(V\in L^\infty(\mathbb R_\tau;L^3(\mathbb R^3_y))\), its Leray
pressure satisfies
\[
   P=R_iR_j(V_iV_j)\in L^\infty(\mathbb R;L^{3/2}(\mathbb R^3)),
\]
and on every compact cylinder
\(K=B_R\times[\tau_1,\tau_2]\) there are functions \(a_n(\tau)\) such that
\[
   P_n-a_n(\tau)
   \rightharpoonup P
   \quad\text{weakly in }L^q(K)
\]
for every \(1<q<\infty\), after passing to a subsequence which may depend on
the compact cylinder and on \(q\).  The functions \(a_n\) may be chosen
measurable on \([\tau_1,\tau_2]\).
\end{theorem}

The second result records the invariance of the \(L^3\)-norm under the
Navier--Stokes scaling and the self-similar change of variables.

\begin{proposition}[Scaling of the \(L^3\)-norm]\label{prop:L3-transfer}
Let \(u\) be a smooth Navier--Stokes solution on
\(\mathbb R^3\times(0,T)\), and let
\[
   u_k(x,t)=\lambda_k u(x_0+\lambda_k x,T+\lambda_k^2t),
   \qquad \lambda_k\downarrow0,
\]
converge locally smoothly on \(\mathbb R^3\times(-\infty,0)\) to an ancient solution
\(U\).  Assume that
\[
   \sup_{0<t<T}\|u(t)\|_{L^3(\mathbb R^3)}\le A<\infty .
\]
Then the self-similar representative
\[
   V(y,\tau)=\sqrt{-t}\,U(y\sqrt{-t},t),
   \qquad \tau=-\log(-t),
\]
satisfies
\[
   \sup_{\tau\in\mathbb R}\|V(\tau)\|_{L^3(\mathbb R^3)}\le A .
\]
\end{proposition}

\begin{theorem}[Escauriaza--Seregin--\v{S}verak critical \(L^3\) criterion]
\label{thm:ESS-L3}
Let \((u,p)\) be a finite-energy suitable weak solution of the
Navier--Stokes equations on \(\mathbb R^3\times(0,T)\).  If
\[
   \sup_{0<t<T}\|u(t)\|_{L^3(\mathbb R^3)}<\infty ,
\]
then no point of \(\mathbb R^3\times\{T\}\) is singular.
\end{theorem}

\begin{proof}
This is the endpoint regularity theorem of Escauriaza--Seregin--\v{S}verak
\cite{ESS2003}.  We use it only through the logically equivalent
contrapositive: if a finite-time singularity occurs at time \(T\), then
\(\|u(t)\|_{L^3(\mathbb R^3)}\) is unbounded as \(t\uparrow T\).
\end{proof}

\begin{corollary}[Exclusion of the source-level critical \(L^3\) regime]
\label{cor:source-L3-branch}
A finite-energy Type I singularity cannot occur for a suitable weak solution
satisfying
\[
   u\in L^\infty(0,T;L^3(\mathbb R^3)).
\]
Consequently the analysis of Type I blow-up limits is relevant only outside
the global critical class \(L^\infty(0,T;L^3(\mathbb R^3))\).
\end{corollary}

\begin{proof}
The asserted \(L^\infty(0,T;L^3(\mathbb R^3))\) bound is exactly the hypothesis of
\Cref{thm:ESS-L3}.  Hence there is no singular point at time \(T\), and in
particular no Type I singular point.  Thus the local concentration and
Type I/Type II alternatives are never reached under this global critical
assumption.
\end{proof}

\begin{theorem}[Small bounded ancient solutions vanish]\label{thm:small}
There exists an absolute constant \(\varepsilon_0>0\) with the following property.  If
\(V\) is a bounded mild ancient solution of \eqref{eq:centered} and
\[
   \|V\|_{L^\infty(\mathbb R^3\times\mathbb R)}\le \varepsilon_0,
\]
then \(V\equiv0\).
\end{theorem}

\begin{theorem}[Stationary limits under uniform \(L^3\)-tightness]\label{thm:stationary-limits}
Let \(V\) be a bounded mild ancient solution of \eqref{eq:centered}.  Assume
that
\begin{equation}\label{eq:L3-bound-tight}
   \sup_{\tau\in\mathbb R}\|V(\tau)\|_{L^3(\mathbb R^3)}<\infty
\end{equation}
and
\begin{equation}\label{eq:tightness}
   \lim_{R\to\infty}
   \sup_{\tau\in\mathbb R}
   \int_{|y|>R}|V(y,\tau)|^3\,dy=0.
\end{equation}
Let \(\tau_n\) be any sequence, and set
\[
   V_n(y,\tau)=V(y,\tau+\tau_n).
\]
Assume that \(V_n\to W\) locally smoothly on \(\mathbb R^3\times\mathbb R\), and that
\(W\) is independent of \(\tau\).  Then \(W\equiv0\).
\end{theorem}

\begin{corollary}[Stationary solutions under uniform \(L^3\)-tightness]\label{cor:stationary-tight}
Let \(V\) be a bounded mild ancient solution such that
\[
   \sup_{\tau\in\mathbb R}\|V(\tau)\|_{L^3(\mathbb R^3)}<\infty
\]
and
\[
   \lim_{R\to\infty}
   \sup_{\tau\in\mathbb R}
   \int_{|y|>R}|V(y,\tau)|^3\,dy=0.
\]
If \(V(y,\tau)=W(y)\), then
\[
   V\equiv0.
\]
\end{corollary}

\begin{corollary}[Positive size threshold]\label{cor:size-threshold}
Every nonzero bounded mild ancient solution of \eqref{eq:centered} satisfies
\[
   \|V\|_{L^\infty(\mathbb R^3\times\mathbb R)}\ge \varepsilon_0,
\]
where \(\varepsilon_0\) is the constant in \cref{thm:small}.
\end{corollary}

\begin{corollary}[Exclusion of the small-amplitude ancient regime]\label{cor:small-branch}
Let \(V\) be a bounded mild ancient solution of \eqref{eq:centered}.  If
\[
   \|V\|_{L^\infty(\mathbb R^3\times\mathbb R)}\le\varepsilon_0,
\]
where \(\varepsilon_0\) is the constant in \cref{thm:small}, then
\[
   V\equiv0.
\]
Consequently no nonzero Type I ancient limit in centered variables can belong
to this small-amplitude class.
\end{corollary}

\begin{proof}
The first assertion is exactly \cref{thm:small}.  The final sentence concerns
only ancient limits obtained from a source-level Type I singular point.  In
that setting the nonvanishing normalization is not an extra assumption of this
paper: it is included in the definition of the source-generated normalized
ancient limit used here.  Therefore a Type I ancient limit satisfying this
normalization cannot both be nonzero in that normalization and vanish by the
small-amplitude theorem.
\end{proof}

\begin{corollary}[Exclusion of the stationary \(L^3\) ancient regime]\label{cor:stationary-L3-branch}
Let \(V\) be a bounded mild ancient solution of \eqref{eq:centered}.  Suppose
that \(V\) is stationary, \(V(y,\tau)=W(y)\), and
\[
   W\in L^3(\mathbb R^3).
\]
Then
\[
   V\equiv0.
\]
Consequently no nonzero Type I ancient limit in centered variables can be both
stationary and globally \(L^3\).
\end{corollary}

\begin{proof}
By \cref{thm:gauges}, \(V\) is smooth.  Since it is independent of \(\tau\),
\(\partial_\tau V=0\) in \(\mathcal D'(\mathbb R^3\times\mathbb R)\).  By
\cref{lem:mild-to-projected}, the mild formulation implies the projected
stationary equation
\[
   0
   =
   \Delta W-\frac12y\cdot\nabla W-\frac12W
   -\mathbb P\operatorname{div}(W\otimes W).
\]
The assumption \(W\in L^3\) implies
\(W_iW_j\in L^{3/2}(\mathbb R^3)\).  Since the Riesz transforms are bounded
on \(L^{3/2}(\mathbb R^3)\) \cite{Stein1970}, the Leray pressure
\[
   P=R_iR_j(W_iW_j)
\]
belongs to \(L^{3/2}(\mathbb R^3)\).  The whole-space Helmholtz decomposition
in \(\mathcal D'(\mathbb R^3)\) \cite{Sohr2001,Galdi2011} gives
\[
   (I-\mathbb P)\operatorname{div}(W\otimes W)=-\nabla P .
\]
Equivalently,
\[
   \mathbb P\operatorname{div}(W\otimes W)
   =
   \operatorname{div}(W\otimes W)+\nabla P .
\]
Substituting this identity into the projected equation gives
\[
   \operatorname{div}(W\otimes W)+\nabla P
   =
   \Delta W-\frac12(W+y\cdot\nabla W)
\]
in \(\mathcal D'(\mathbb R^3)\).  Since \(\operatorname{div}W=0\),
\[
   \operatorname{div}(W\otimes W)_i
   =
   \partial_j(W_iW_j)
   =
   W_j\partial_jW_i+W_i\partial_jW_j
   =
   (W\cdot\nabla)W_i .
\]
Thus \eqref{eq:stationary-profile} holds in \(\mathcal D'(\mathbb R^3)\), and
\cref{thm:NRS} gives \(W\equiv0\).  The final claim is then a direct
comparison with the nonvanishing normalization supplied by the Type I blow-up
normalization imposed on source-generated Type I ancient limits.  No
additional normalization or compactness assumption is introduced in this
corollary.
\end{proof}

\section{Pressure estimates and compactness}

We first record the local pressure estimates used below.  Time-dependent
pressure gauges are harmless for source suitable weak solutions and their
parabolic rescalings because they have zero spatial gradient and do not affect
the local energy inequality.  The lemmas in this section are the whole-space
Leray-pressure and smooth-ancient-profile
estimates needed for the conditional arguments of the present paper.

\begin{lemma}[Local pressure decomposition]\label{lem:pressure-decomp}
Let \(V\) be bounded on \(B_{4R}\times[\tau_1,\tau_2]\), and let
\(P\in L^1(B_{4R}\times[\tau_1,\tau_2])\) solve
\[
   -\Delta P=\partial_i\partial_j(V_iV_j)
\]
in \(\mathcal D'(B_{4R}\times(\tau_1,\tau_2))\).  Then there is a measurable
function \(a:(\tau_1,\tau_2)\to\mathbb R\) such that,
for every \(1<q<\infty\) and almost every \(\tau\in[\tau_1,\tau_2]\),
\[
   \|P(\tau)-a(\tau)\|_{L^q(B_R)}
   \le C(q,R)\|V(\tau)\|_{L^\infty(B_{4R})}^2
   +C(q,R)\|P(\tau)\|_{L^1(B_{4R})}.
\]
If \(P\) is the whole-space Leray pressure and
\(V\in L^\infty(\mathbb R_\tau;L^3(\mathbb R^3_y))\),
then, after choosing the Leray gauge,
\[
   \|P-a(\tau)\|_{L^q(B_R\times[\tau_1,\tau_2])}
   \le C(q,R,\tau_1,\tau_2)
      \left(
      \|V\|_{L^\infty(B_{4R}\times[\tau_1,\tau_2])}^2
      +\sup_{\tau\in\mathbb R}\|V(\tau)\|_{L^3(\mathbb R^3)}^2
      \right).
\]
Here the spacetime norm is taken with the representative
\((y,\tau)\mapsto P(y,\tau)-a(\tau)\).
\end{lemma}

\begin{proof}
For the first assertion it is enough to argue for almost every fixed
\(\tau\), with constants independent of that \(\tau\).  Indeed, testing the
spacetime distributional equation against products
\(\zeta(\tau)\varphi(y)\) and applying Fubini shows that, for every
\(\varphi\in C_c^\infty(B_{4R})\), the scalar identity
\[
   \int_{B_{4R}}P(y,\tau)(-\Delta\varphi)(y)\,dy
   =
   \int_{B_{4R}}V_i(y,\tau)V_j(y,\tau)\partial_i\partial_j\varphi(y)\,dy
\]
holds for almost every \(\tau\).  Applying this first to a countable dense
set of test functions in \(C_c^\infty(B_{4R})\), and then using continuity of
the two sides with respect to the usual test-function topology, gives a
single null set outside which
\[
   -\Delta P(\cdot,\tau)
   =
   \partial_i\partial_j(V_iV_j)(\cdot,\tau)
   \quad\text{in }\mathcal D'(B_{4R})
\]
for every test function.  Choose
\(\chi\in C_c^\infty(B_{4R})\) with \(\chi\equiv1\) on \(B_{3R}\), and set
\[
   \pi_{\mathrm{loc}}
   =
   (-\Delta)^{-1}\partial_i\partial_j(\chi V_iV_j).
\]
The inverse Laplacian is the whole-space Newtonian potential, so
\(\pi_{\mathrm{loc}}\) is defined as a singular integral of the compactly
supported function
\(\chi V_iV_j\).  Since \(V\) is bounded on \(B_{4R}\), this compactly
supported product belongs to \(L^q(\mathbb R^3)\) for every
\(1<q<\infty\).  The Calder\'on--Zygmund theorem \cite{Stein1970} gives
\[
   \|\pi_{\mathrm{loc}}\|_{L^q(\mathbb R^3)}
   \le C(q)\|\chi V_iV_j\|_{L^q(\mathbb R^3)}
   \le C(q,R)\|V\|_{L^\infty(B_{4R})}^2 .
\]
In \(B_{3R}\) we have
\[
   -\Delta(P-\pi_{\mathrm{loc}})
   =
   \partial_i\partial_j(V_iV_j)-\partial_i\partial_j(\chi V_iV_j)=0
\]
in \(\mathcal D'(B_{3R})\), because \(\chi\equiv1\) on \(B_{3R}\).  Weyl's lemma then
implies that \(h=P-\pi_{\mathrm{loc}}\) is harmonic in \(B_{3R}\); this is
the interior elliptic regularity implication for the Laplacian, see for
instance \cite{GilbargTrudinger2001}.  Define
\[
   a(\tau)=(h)_{B_{2R}}(\tau)
   =|B_{2R}|^{-1}\int_{B_{2R}}h(y,\tau)\,dy .
\]
The mean value property and interior estimates for harmonic functions give
\[
   \|h-a(\tau)\|_{L^q(B_R)}
   \le C(R,q)\|h-a(\tau)\|_{L^1(B_{2R})}.
\]
Since \(a(\tau)\) is the mean of \(h\) on \(B_{2R}\),
\[
   \|h-a(\tau)\|_{L^1(B_{2R})}
   \le
   \|h\|_{L^1(B_{2R})}
   +|B_{2R}|\, |(h)_{B_{2R}}|
   \le 2\|h\|_{L^1(B_{2R})}.
\]
Therefore
\[
   \|h-a(\tau)\|_{L^q(B_R)}
   \le C(R,q)\|h\|_{L^1(B_{2R})}.
\]
Since \(h=P-\pi_{\mathrm{loc}}\), the last two estimates imply
\[
   \|P-a(\tau)\|_{L^q(B_R)}
   \le C(R,q)\|P\|_{L^1(B_{2R})}
      +C(R,q)\|\pi_{\mathrm{loc}}\|_{L^1(B_{2R})}
      +C(R,q)\|\pi_{\mathrm{loc}}\|_{L^q(B_R)} .
\]
The \(\pi_{\mathrm{loc}}\)-terms are bounded by the Calder\'on--Zygmund
estimate above, and this gives the pointwise-in-time local pressure estimate.
The function \(a\) is measurable because it is obtained by integrating the
locally integrable function \(h=P-\pi_{\mathrm{loc}}\) over \(B_{2R}\).

Assume now that \(P=R_iR_j(V_iV_j)\) is the Leray pressure and
\[
   M_3=\sup_{\tau\in\mathbb R}\|V(\tau)\|_{L^3(\mathbb R^3)}<\infty .
\]
We spell out the gauge because this is where the nonlocal pressure tail
enters.  Let \(K_{ij}\) be the homogeneous
degree \(-3\) kernel of \(R_iR_j\), interpreted in the principal value sense.
The singular-integral facts used here are the Riesz-transform
Calder\'on--Zygmund bounds \cite{Stein1970}.  The local part
\(R_iR_j(\chi V_iV_j)\) is controlled exactly as above.  For the exterior part
we subtract a constant in \(x\).  The kernel is smooth away from the origin
and satisfies, for
\(x\in B_R\) and \(|z|>4R\),
\[
   |K_{ij}(x-z)-K_{ij}(-z)|\le CR|z|^{-4}.
\]
On the dyadic annulus
\[
   A_m=\{2^mR<|z|<2^{m+1}R\},\qquad m\ge2,
\]
H\"older's inequality gives
\[
   \int_{A_m}|V(z,\tau)|^2\,dz
   \le
   \|V(\tau)\|_{L^3(A_m)}^2 |A_m|^{1/3}
   \le C M_3^2 2^mR .
\]
Consequently
\[
   \sum_{m=2}^\infty
   R(2^mR)^{-4}
   \int_{2^mR<|z|<2^{m+1}R}|V(z,\tau)|^2\,dz
   \le
   C M_3^2 R^{-2},
\]
uniformly in \(\tau\).  This gives a bounded far-field contribution modulo
the chosen constant, with a bound depending only on \(M_3\).  Combining the
near-field Calder\'on--Zygmund estimate and this dyadic exterior estimate gives
the stated Leray-gauge bound on \(B_R\times[\tau_1,\tau_2]\).
\end{proof}

\begin{lemma}[Smoothing of bounded mild solutions]\label{lem:mild-smoothing}
Let \(V\) be a bounded mild solution of \eqref{eq:centered} on
\(\mathbb R^3\times(\tau_0,\tau_1)\).  For every compact cylinder
\(K\Subset\mathbb R^3\times(\tau_0,\tau_1)\) and every integer \(m\ge0\),
\[
   \|V\|_{C^m(K)}
   \le
   C\bigl(m,K,\tau_0,\tau_1,\|V\|_{L^\infty}\bigr).
\]
The same bound holds uniformly for families with a common \(L^\infty\) bound.
\end{lemma}

\begin{proof}
Fix a compact interval \(J\Subset(\tau_0,\tau_1)\) containing the time
projection of \(K\), and put \(t=-e^{-\tau}\).  Let
\(t(J)=\{-e^{-\tau}:\tau\in J\}\).  On this compact negative time interval
define
\[
   U(x,t)=(-t)^{-1/2}V(x/\sqrt{-t},-\log(-t)).
\]
The map
\[
   (x,t)\longmapsto (y,\tau)=\left(x/\sqrt{-t},-\log(-t)\right)
\]
is a \(C^\infty\) diffeomorphism from
\(\mathbb R^3\times t(J)\) onto \(\mathbb R^3\times J\).  Its spatial
Jacobian at fixed time is \((-t)^{-3/2}\), and all derivatives of the map and
of its inverse are bounded on compact subcylinders because \(t(J)\Subset
(-\infty,0)\).
We first justify that \(U\) is a bounded mild Navier--Stokes solution on
\(\mathbb R^3\times t(J)\).  If
\[
   x=\sqrt{-t}\,y,\qquad t=-e^{-\tau},
\]
then \(d\tau/dt=-1/t\) and
\[
   \partial_t U
   =
   (-t)^{-3/2}
   \left(
      \partial_\tau V
      +\frac12V+\frac12y\cdot\nabla_yV
   \right),
   \quad
   \Delta_xU=(-t)^{-3/2}\Delta_yV,
\]
while
\[
   (U\cdot\nabla_x)U=(-t)^{-3/2}(V\cdot\nabla_y)V .
\]
The pressure transforms as
\[
   p(x,t)=(-t)^{-1}P(x/\sqrt{-t},-\log(-t)).
\]
Consequently
\[
   \nabla_xp=(-t)^{-3/2}\nabla_yP,
   \qquad
   \operatorname{div}_xU=(-t)^{-1}\operatorname{div}_yV .
\]
Substituting these identities in \eqref{eq:centered} gives
\[
   \partial_tU+(U\cdot\nabla)U+\nabla p=\Delta U,
   \qquad \operatorname{div}U=0,
\]
in \(\mathcal D'(\mathbb R^3\times t(J))\).  This distributional statement
is obtained by testing against compactly supported functions in
\((x,t)\), pulling the tests back by the above diffeomorphism, and using the
displayed Jacobian; no boundary term occurs because the tests are compactly
supported in the open time interval \(t(J)\).  Moreover, for
\(t_1<t_2\) in \(t(J)\), the corresponding logarithmic times
\(\tau_i=-\log(-t_i)\) satisfy \(\tau_1<\tau_2\), and testing the weak-*
Duhamel formula \eqref{eq:mild} against the rescaled function
\((-t_2)^{-3/2}\phi(\cdot/\sqrt{-t_2})\) gives, after the change of
variables \(x=\sqrt{-t}\,y\), the usual heat-kernel Duhamel formula
\[
   U(t_2)=e^{(t_2-t_1)\Delta}U(t_1)
   -\int_{t_1}^{t_2}
      e^{(t_2-s)\Delta}\mathbb P\operatorname{div}(U\otimes U)(s)\,ds
\]
in weak-* \(L^\infty(\mathbb R^3)\).  The covariance of the Leray projection
under translations and dilations is used here through its Fourier multiplier
symbol
\(\delta_{ij}-\xi_i\xi_j/|\xi|^2\), which is homogeneous of degree zero.
Thus \(U\) is a bounded mild Navier--Stokes solution on every compact
subinterval of \(t(J)\).  On \(J\), the quantities \((-t)^{1/2}\) and
\((-t)^{-1/2}\) are bounded above and below by positive constants depending
only on \(J\).  Hence
\[
   \|U\|_{L^\infty(\mathbb R^3\times t(J))}
   \le
   C(J)\|V\|_{L^\infty(\mathbb R^3\times(\tau_0,\tau_1))}.
\]
The interior smoothing theorem for bounded mild Navier--Stokes solutions,
in the form of Kato \cite{Kato1984} and Giga--Miyakawa
\cite{GigaMiyakawa1985}, followed by the local parabolic bootstrap
for the Stokes system \cite{Ladyzhenskaya1968,Sohr2001}, gives bounds for all
spatial and time derivatives of \(U\) on compact subcylinders of
\(\mathbb R^3\times t(J)\).  The constants depend only on the displayed
bound, the compact cylinder, and the parabolic distance to the ends of
\(t(J)\).  The derivatives of the map
\((y,\tau)\mapsto(x,t)=(y\sqrt{-t},-e^{-\tau})\) are bounded on the chosen
compact cylinder, with bounds depending only on \(J\) and the spatial size of
the cylinder.  Returning to \((y,\tau)\) therefore gives the asserted
estimates for \(V\).  The same argument applies uniformly to any family with
a common \(L^\infty\) bound.
\end{proof}

\begin{lemma}[From mild to projected distributional form]\label{lem:mild-to-projected}
Let \(V\) be a bounded mild solution of \eqref{eq:centered} on
\(\mathbb R^3\times(\tau_0,\tau_1)\).  If \(V\) is smooth on this cylinder,
then \(V\) satisfies \eqref{eq:projected} in
\(\mathcal D'(\mathbb R^3\times(\tau_0,\tau_1))\).
\end{lemma}

\begin{proof}
Write
\[
   L=\Delta-\frac12y\cdot\nabla-\frac12,
   \qquad
   N(\tau)=\mathbb P\operatorname{div}(V\otimes V)(\tau).
\]
The semigroup \(S(a)\) is generated by \(L\) on test functions and, by
duality, on the \(L^\infty\)-weak-* side.  This is the differentiation
theorem for \(C_0\)-semigroups applied to compactly supported smooth tests
\cite{Pazy1983}.  Fix
\(\phi\in C_c^\infty(\mathbb R^3;\mathbb R^3)\) and
\(\tau\in(\tau_0,\tau_1)\).  For \(h>0\) small enough that
\(\tau+h<\tau_1\), the mild formula on \([\tau,\tau+h]\) gives
\[
   V(\tau+h)-V(\tau)
   =
   (S(h)-I)V(\tau)
   -
   \int_0^h S(h-s)N(\tau+s)\,ds .
\]
Pair this identity with \(\phi\) and divide by \(h\).  For the linear term,
\[
   \frac1h\langle (S(h)-I)V(\tau),\phi\rangle
   =
   \left\langle V(\tau),\frac{S(h)^*\phi-\phi}{h}\right\rangle
   \to
   \langle V(\tau),L^*\phi\rangle
   =
   \langle LV(\tau),\phi\rangle ,
\]
where the last equality follows by integration by parts, justified because
\(\phi\) is compactly supported and \(V(\tau)\) is smooth.  For the Duhamel
term put \(F=V\otimes V\).  By definition of \(N\),
\[
   \langle S(h-s)N(\tau+s),\phi\rangle
   =
   \langle F(\tau+s),
      (S(h-s)\mathbb P\operatorname{div})^*\phi\rangle .
\]
For fixed \(\phi\in C_c^\infty\), the Oseen-kernel representation and the
Riesz-transform bounds imply
\[
   (S(r)\mathbb P\operatorname{div})^*\phi
   \to
   (\mathbb P\operatorname{div})^*\phi
   \quad\text{in }L^1(\mathbb R^3)
   \qquad (r\downarrow0),
\]
and the same family is bounded in \(L^1\) for \(0\le r\le h_0\), with \(h_0\)
fixed.  The convergence in \(L^1\) also implies uniform tightness of this
family in \(L^1\) for small \(r\).  This is the strong continuity of the
regularized Stokes kernel on smooth test data \cite{Stein1970,Sohr2001}.
To see explicitly why the limiting adjoint is integrable, use that
\(\mathbb P\) is self-adjoint on \(L^2\) and commutes with derivatives.  For
matrix fields \(F\),
\[
   \langle \mathbb P\operatorname{div}F,\phi\rangle
   =
   -\int_{\mathbb R^3}F_{ij}\,
      \partial_j(\mathbb P\phi)_i\,dy .
\]
Hence
\[
   (\mathbb P\operatorname{div})^*\phi
   =
   -\nabla(\mathbb P\phi).
\]
Since \(\phi\) is compactly supported and smooth, the local part is smooth and
compactly supported, while the nonlocal Riesz part has kernel derivatives
decaying like \(|y|^{-4}\).  This decay is integrable in three dimensions,
so \((\mathbb P\operatorname{div})^*\phi\in L^1(\mathbb R^3)\).
Since \(F\) is bounded and smooth, \(F(\tau+s)\to F(\tau)\) locally uniformly
and pointwise as \(s\to0\).  We spell out the limiting argument.  Let
\(\varepsilon>0\).  By the uniform \(L^1\)-tightness of
\((S(r)\mathbb P\operatorname{div})^*\phi\), choose \(R\) so large that the
\(L^1(\mathbb R^3\setminus B_R)\)-norm of every member of the family, and of
the limit \((\mathbb P\operatorname{div})^*\phi\), is at most
\(\varepsilon/(1+\|F\|_{L^\infty})\).  On \(B_R\), local uniform convergence
of \(F(\tau+s)\) and strong \(L^1\)-convergence of the adjoint kernels give
convergence of the pairings.  The exterior part is bounded by
\(C\|F\|_{L^\infty}\varepsilon\), uniformly for small \(s\) and \(r\).  Hence
the integrand converges to
\(\langle N(\tau),\phi\rangle\), and it is bounded by an integrable constant
on \(0<s<h\) for \(h<h_0\).  Averaging over \(s\in(0,h)\) therefore preserves
the limit as \(h\downarrow0\).
Therefore
\[
   \frac1h\int_0^h
      \langle N(\tau+s),S(h-s)^*\phi\rangle\,ds
   \to
   \langle N(\tau),\phi\rangle .
\]
The left-hand side converges to
\(\langle\partial_\tau V(\tau),\phi\rangle\) because \(V\) is smooth in
\(\tau\).  Thus
\[
   \langle\partial_\tau V(\tau),\phi\rangle
   =
   \langle LV(\tau),\phi\rangle-\langle N(\tau),\phi\rangle
\]
for every compactly supported smooth \(\phi\) and every
\(\tau\in(\tau_0,\tau_1)\).  Multiplying by a scalar test function of time and
integrating in \(\tau\) gives \eqref{eq:projected} in
\(\mathcal D'(\mathbb R^3\times(\tau_0,\tau_1))\).
\end{proof}

\begin{lemma}[Local smooth compactness]\label{lem:smooth-compactness}
Let \(V_n\) be bounded mild ancient solutions of \eqref{eq:centered} with
\[
   \sup_n\|V_n\|_{L^\infty(\mathbb R^3\times\mathbb R)}\le M.
\]
Then, after passing to a subsequence, \(V_n\) converges in
\(C^m_{\mathrm{loc}}(\mathbb R^3\times\mathbb R)\) for every \(m\ge0\) to a bounded mild ancient
solution \(V\).
\end{lemma}

\begin{proof}
\Cref{lem:mild-smoothing} gives uniform \(C^m\)-bounds on every compact
cylinder, for every \(m\).  Choose an exhaustion
\[
   Q_\ell=B_\ell\times[-\ell,\ell],\qquad \ell=1,2,\dots .
\]
For fixed \(\ell\) and \(m\), the uniform \(C^{m+1}(Q_{\ell+1})\)-bound gives
equicontinuity of all derivatives of order at most \(m\) on \(Q_\ell\).
Arzel\`a--Ascoli gives a subsequence converging in \(C^m(Q_\ell)\).  Extracting
successively over the countable pairs \((\ell,m)\) and taking the diagonal
subsequence yields convergence in \(C^m_{\mathrm{loc}}\) for every \(m\).
The limit satisfies
\[
   \|V\|_{L^\infty(\mathbb R^3\times\mathbb R)}\le M
\]
by pointwise convergence on compact subsets.  It is also divergence free in
the sense of distributions: for every
\(\varphi\in C_c^\infty(\mathbb R^3\times\mathbb R)\),
\[
   \langle \operatorname{div}V,\varphi\rangle
   =
   -\int V\cdot\nabla\varphi
   =
   -\lim_{n\to\infty}\int V_n\cdot\nabla\varphi
   =
   0,
\]
because \(\operatorname{div}V_n=0\) distributionally.

It remains to verify that the global mild formulation is closed under this
convergence.  Fix \(\sigma<\tau\).  For every \(\phi\in L^1(\mathbb R^3)\), local
uniform convergence, the common \(L^\infty\) bound, and approximation of
\(\phi\) by compactly supported functions give
\[
   \int_{\mathbb R^3} V_n(\cdot,t)\phi\,dy
   \to
   \int_{\mathbb R^3} V(\cdot,t)\phi\,dy,
   \qquad t\in\{\sigma,\tau\}.
\]
Indeed, if \(\phi=\phi_R+(\phi-\phi_R)\) with \(\phi_R\in C_c^\infty\) and
\(\|\phi-\phi_R\|_{L^1}\to0\), then the compactly supported part converges by
uniform convergence on a compact set, while the remainder is bounded by
\(2M\|\phi-\phi_R\|_{L^1}\).  The same argument applies to
\(V_n\otimes V_n\) against \(L^1\) matrix-valued test functions, with
\(2M^2\) in place of \(2M\).

Now test \eqref{eq:mild} against a fixed
\(\phi\in C_c^\infty(\mathbb R^3;\mathbb R^3)\).  By duality from
\cref{lem:stokes-estimates}, the adjoint of
\(S(a)\mathbb P\operatorname{div}\) satisfies
\[
   \|(S(a)\mathbb P\operatorname{div})^*\phi\|_{L^1}
   \le
   C e^{-a/2}(1-e^{-a})^{-1/2}\|\phi\|_{L^1}.
\]
For each \(s\in(\sigma,\tau)\) the preceding paragraph gives
\[
   \langle V_n\otimes V_n(s), (S(\tau-s)\mathbb P\operatorname{div})^*\phi\rangle
   \to
   \langle V\otimes V(s), (S(\tau-s)\mathbb P\operatorname{div})^*\phi\rangle .
\]
The absolute value of the integrand is bounded by
\[
   C M^2 e^{-(\tau-s)/2}(1-e^{-(\tau-s)})^{-1/2}\|\phi\|_{L^1}.
\]
This function is integrable on \((\sigma,\tau)\): near \(s=\tau\) it is
comparable to \((\tau-s)^{-1/2}\), and away from \(s=\tau\) it is bounded.
Dominated convergence therefore passes the Duhamel term to the limit.  The
linear terms at times \(\sigma\) and \(\tau\) pass to the limit by the first
paragraph.  Thus \eqref{eq:mild} holds for \(V\) against every compactly
supported smooth test function.  Since both sides are bounded functions of
space, density of \(C_c^\infty(\mathbb R^3)\) in \(L^1(\mathbb R^3)\) extends
the identity to all \(L^1\) tests, which is exactly weak-* equality in
\(L^\infty(\mathbb R^3)\).
\end{proof}

\begin{proof}[Proof of \cref{thm:gauges}]
Smoothness follows directly from \cref{lem:mild-smoothing}.  For each fixed
\(\tau\), local uniform convergence gives
\[
   \int_{B_R}|V(y,\tau)|^3\,dy
   =
   \lim_{n\to\infty}\int_{B_R}|V_n(y,\tau)|^3\,dy
\]
for every \(R<\infty\).  Letting \(R\to\infty\) shows that
\(V\in L^\infty(\mathbb R_\tau;L^3(\mathbb R^3_y))\).  More explicitly, if
\[
   M_3=\sup_n\sup_{\tau\in\mathbb R}\|V_n(\tau)\|_{L^3},
\]
then monotone convergence in \(R\) yields
\[
   \|V(\tau)\|_{L^3}^3
   =
   \sup_{R<\infty}\int_{B_R}|V(y,\tau)|^3\,dy
   \le M_3^3
\]
for every \(\tau\).  The same smoothing lemma and the uniqueness of the
locally uniform limit imply, after passing to a subsequence if necessary, that
\(V_n\to V\) locally smoothly.  Indeed, any subsequence of \(\{V_n\}\) has a
further subsequence converging in \(C^m_{\mathrm{loc}}\), for every \(m\), by
\cref{lem:smooth-compactness}; its \(C^0_{\mathrm{loc}}\)-limit must be the
given locally uniform limit \(V\).  Thus the subsequence used below may be
chosen so that the convergence is locally smooth.
Since \(V_iV_j\in L^\infty(\mathbb R;L^{3/2}(\mathbb R^3))\), the
Calder\'on--Zygmund boundedness of the Riesz transforms on \(L^{3/2}\)
\cite{Stein1970} defines
\[
   P=R_iR_j(V_iV_j)
   \in L^\infty(\mathbb R;L^{3/2}(\mathbb R^3)).
\]

We now identify the pressure limit in the Leray gauge.  Fix
\(B_R\times I\).  Let \(K_{ij}\) be the kernel of \(R_iR_j\), and choose
\(\eta_R\in C^\infty(\mathbb R^3)\) with \(\eta_R=0\) on \(B_{2R}\) and
\(\eta_R=1\) on \(\mathbb R^3\setminus B_{4R}\).  Set
\[
   b_n(\tau)=\int_{\mathbb R^3}K_{ij}(-z)\eta_R(z)
             V_{n,i}(z,\tau)V_{n,j}(z,\tau)\,dz
\]
and define \(b(\tau)\) analogously with \(V\).  These integrals converge
absolutely and uniformly in \(n,\tau\), since on each dyadic annulus
\(A_\rho=\{\rho<|z|<2\rho\}\)
\[
   \int_{A_\rho}|z|^{-3}|V_n(z,\tau)|^2\,dz
   \le
   C\rho^{-3}\|V_n(\tau)\|_{L^3(A_\rho)}^2 |A_\rho|^{1/3}
   \le C M_3^2\rho^{-2}.
\]
The integrands are measurable in \(\tau\), and the preceding summable
majorant is independent of \(\tau\); hence \(b_n\) and \(b\) are measurable
functions on \(I\).  Consequently the gauges
\[
   a_n(\tau)=b_n(\tau)-b(\tau)
\]
are measurable.
For \(x\in B_R\), the gauged pressure is represented by
\[
   P_n(x,\tau)-b_n(\tau)
   =
   \operatorname{p.v.}\!\int_{\mathbb R^3}
   \bigl(K_{ij}(x-z)-\eta_R(z)K_{ij}(-z)\bigr)
   V_{n,i}(z,\tau)V_{n,j}(z,\tau)\,dz .
\]
This is the principal-value definition of \(R_iR_j\), with the subtracted
term \(b_n(\tau)\) removing only a spatially constant far-field piece;
subtracting such a constant does not change the pressure gradient.

Fix \(L>4R\).  Choose \(\chi_L\in C_c^\infty(B_{2L})\) with
\(\chi_L\equiv1\) on \(B_L\).  Split the gauged pressure into the part with
factor \(\chi_L(z)\) and the remaining exterior part.  On the bounded set
\(B_{2L}\times I\), local smooth convergence implies
\[
   V_{n,i}V_{n,j}\to V_iV_j
   \quad\text{strongly in }L^q(B_{2L}\times I)
\]
for every finite \(q\).  The truncated principal value operator is a
Calder\'on--Zygmund operator applied to \(\chi_L V_{n,i}V_{n,j}\), followed
by restriction to \(B_R\), plus the smooth subtraction operator with kernel
\(\eta_R(z)K_{ij}(-z)\chi_L(z)\).  This subtraction kernel is smooth and
compactly supported because \(\chi_L\) is compactly supported and
\(\eta_R=0\) on \(B_{2R}\), away from the singularity of \(K_{ij}(-z)\) at
the origin.  Calder\'on--Zygmund boundedness on \(L^q\),
\(1<q<\infty\), and Young's inequality for this smooth compactly supported
kernel therefore show that the contribution with factor \(\chi_L\)
converges strongly in \(L^q(B_R\times I)\) for every \(1<q<\infty\).

The exterior contribution is uniform in \(n\).  On the support of
\(1-\chi_L\) one has \(|z|>L\), and for \(x\in B_R\)
\[
   |K_{ij}(x-z)-K_{ij}(-z)|\le C_R|z|^{-4},
\]
and therefore
\[
\begin{aligned}
   &\sup_{n,\tau}\sup_{x\in B_R}
      \left|
         \int_{|z|>L}[K_{ij}(x-z)-K_{ij}(-z)]
         V_{n,i}V_{n,j}(z,\tau)\,dz
      \right| \\
   &\le
   C_R\sum_{m=0}^\infty (2^mL)^{-4}
      \int_{2^mL<|z|<2^{m+1}L}|V_n(z,\tau)|^2\,dz  \\
   &\le C_R M_3^2 L^{-3}.
\end{aligned}
\]
The same estimate holds for the limit \(V\).  Letting \(n\to\infty\) for fixed
\(L\), and then \(L\to\infty\), gives
\(P_n-b_n\to P-b\) strongly in \(L^q(B_R\times I)\) for every
\(1<q<\infty\).  Taking \(a_n=b_n-b\) gives the stated convergence to the
Leray pressure \(P\); in particular the asserted weak convergence follows.
\end{proof}

\section{\texorpdfstring{\(L^3\)}{L3} transfer and tightness}

The \(L^3\)-norm is invariant under both the Navier--Stokes parabolic scaling
and the centered self-similar scaling.  The calculation below records the
resulting global \(L^3\) consequence used here.

\begin{proof}[Proof of \cref{prop:L3-transfer}]
For each fixed \(t<0\), the time \(T+\lambda_k^2t\) belongs to \((0,T)\) for
all sufficiently large \(k\).  For those \(k\),
\[
   \|u_k(t)\|_{L^3(\mathbb R^3)}
   =
   \|u(T+\lambda_k^2t)\|_{L^3(\mathbb R^3)}
   \le A.
\]
The equality follows by the change of variables
\(X=x_0+\lambda_kx\), for which \(dX=\lambda_k^3dx\) and
\(|u_k(x,t)|^3=\lambda_k^3|u(X,T+\lambda_k^2t)|^3\).  By local smooth
convergence, \(u_k(\cdot,t)\to U(\cdot,t)\) uniformly on each ball \(B_R\).
Since \(B_R\) has finite measure and
\[
   \|u_k(\cdot,t)-U(\cdot,t)\|_{L^3(B_R)}
   \le |B_R|^{1/3}
      \|u_k(\cdot,t)-U(\cdot,t)\|_{L^\infty(B_R)},
\]
the convergence is strong in \(L^3(B_R)\).  Hence
\[
   \int_{B_R}|U(x,t)|^3\,dx
   =
   \lim_{k\to\infty}\int_{B_R}|u_k(x,t)|^3\,dx
   \le A^3
\]
for every \(R<\infty\) and every \(t<0\).  Since
\(\int_{B_R}|U(x,t)|^3\,dx\) is monotone increasing in \(R\), the monotone
convergence theorem gives
\(\|U(t)\|_{L^3(\mathbb R^3)}\le A\) for every \(t<0\).
Finally,
\[
   \|V(\tau)\|_{L^3(\mathbb R^3)}
   =
   \|U(t)\|_{L^3(\mathbb R^3)},
   \qquad t=-e^{-\tau},
\]
because, with \(x=y\sqrt{-t}\),
\[
   \int_{\mathbb R^3}|V(y,\tau)|^3\,dy
   =
   \int_{\mathbb R^3}(-t)^{3/2}|U(y\sqrt{-t},t)|^3\,dy
   =
   \int_{\mathbb R^3}|U(x,t)|^3\,dx .
\]
\end{proof}

\begin{lemma}[Closure of uniform \(L^3\)-tightness]\label{lem:tight-closure}
Let \(V_n\) be bounded mild ancient solutions of \eqref{eq:centered}.  Assume
that
\[
   \sup_n\|V_n\|_{L^\infty(\mathbb R^3_y\times\mathbb R_\tau)}
   +\sup_n\sup_{\tau\in\mathbb R}\|V_n(\tau)\|_{L^3(\mathbb R^3_y)}
   <\infty
\]
and that the tightness in \eqref{eq:tightness} is uniform in \(n\):
\[
   \lim_{R\to\infty}
   \sup_n\sup_{\tau\in\mathbb R}\int_{|y|>R}|V_n(y,\tau)|^3\,dy=0.
\]
If \(V_n\to V\) locally smoothly, then
\[
   \sup_{\tau\in\mathbb R}\|V(\tau)\|_{L^3(\mathbb R^3)}<\infty
\]
and
\[
   \lim_{R\to\infty}
   \sup_{\tau\in\mathbb R}\int_{|y|>R}|V(y,\tau)|^3\,dy=0.
\]
\end{lemma}

\begin{proof}
Let
\[
   M_3=\sup_n\sup_{\tau\in\mathbb R}\|V_n(\tau)\|_{L^3}<\infty .
\]
Fix \(\tau\).  On \(B_R\), local smooth convergence gives
\[
   \int_{B_R}|V(\tau)|^3\,dy
   =
   \lim_{n\to\infty}\int_{B_R}|V_n(\tau)|^3\,dy.
\]
Taking the supremum over \(R\) and using monotone convergence gives
\[
   \|V(\tau)\|_{L^3}^3
   =
   \sup_{R<\infty}\int_{B_R}|V(\tau)|^3\,dy
   \le M_3^3 .
\]
Thus \(V\) has the same uniform \(L^3\) bound.  For tightness, choose \(R\)
large so that
\[
   \sup_n\sup_{\tau\in\mathbb R}\int_{|y|>R}|V_n(y,\tau)|^3\,dy
   \le \varepsilon .
\]
Local smooth convergence gives pointwise convergence everywhere after fixing
\(\tau\).  Fatou's lemma on the exterior domain gives
\[
   \int_{|y|>R}|V(\tau)|^3\,dy
   \le
   \liminf_{n\to\infty}\int_{|y|>R}|V_n(\tau)|^3\,dy
   \le \varepsilon .
\]
The choice of \(R\) was made only from the uniform tail bound for the
sequence \(\{V_n\}\), and therefore is independent of the fixed time
\(\tau\).  Hence
\[
   \sup_{\tau\in\mathbb R}\int_{|y|>R}|V(y,\tau)|^3\,dy
   \le \varepsilon .
\]
Since \(\varepsilon>0\) was arbitrary, the displayed uniform tail limit is
zero.
\end{proof}

\section{The small ancient Liouville theorem}

The proof of \cref{thm:small} uses only the mild formulation and the linear
decay of the centered Stokes flow.  No spatial tightness is needed.

\begin{lemma}[Centered Stokes estimates]\label{lem:stokes-estimates}
There is an absolute constant \(C_0\) such that for every \(a>0\),
\[
   \|S(a)f\|_{L^\infty}
   \le e^{-a/2}\|f\|_{L^\infty},
\]
and
\begin{equation}\label{eq:nonlinear-semigroup-est}
   \|S(a)\mathbb P\operatorname{div} F\|_{L^\infty}
   \le
   C_0 e^{-a/2}(1-e^{-a})^{-1/2}\|F\|_{L^\infty}.
\end{equation}
Consequently
\begin{equation}\label{eq:semigroup-integral}
   C_*:=
   C_0\int_0^\infty e^{-a/2}(1-e^{-a})^{-1/2}\,da
   <\infty .
\end{equation}
\end{lemma}

\begin{proof}
The first estimate follows immediately from the positive kernel formula
\eqref{eq:OU-semigroup}, since the Gaussian factor has total mass one.  The
Oseen-kernel estimates used below are the classical heat-kernel bounds for
the Stokes semigroup; see, for example, \cite{Sohr2001,Galdi2011}.

For the second estimate, first take \(F\in C_c^\infty(\mathbb R^3;\mathbb R^{3\times3})\).
The scalar Ornstein--Uhlenbeck semigroup in \eqref{eq:OU-semigroup} may be
written as
\[
   S(a)g(y)=e^{-a/2}\bigl(e^{(1-e^{-a})\Delta}g\bigr)(e^{-a/2}y).
\]
The Leray projection is the Fourier multiplier
\[
   \widehat{(\mathbb P f)}_i(\xi)
   =
   \left(\delta_{ij}-\frac{\xi_i\xi_j}{|\xi|^2}\right)\widehat f_j(\xi),
   \qquad \xi\ne0.
\]
It is homogeneous of degree zero and therefore commutes with the heat
multiplier \(e^{-s|\xi|^2}\) and with the derivative multipliers
\(i\xi_k\).  Therefore, for
\(s=1-e^{-a}\),
\[
   S(a)\mathbb P\operatorname{div}F(y)
   =
   e^{-a/2}
   \bigl(e^{s\Delta}\mathbb P\operatorname{div}F\bigr)(e^{-a/2}y).
\]
The operator \(e^{s\Delta}\mathbb P\operatorname{div}\) has the Oseen
derivative kernel
\[
   K_s(x)=s^{-2}K_1(x/\sqrt{s}),
\]
with
\[
   \|K_s\|_{L^1(\mathbb R^3)}\le C s^{-1/2}.
\]
This scaling follows from the single derivative in \(\operatorname{div}\):
the heat kernel contributes \(s^{-3/2}\), the derivative contributes
\(s^{-1/2}\), and the change of variables in the \(L^1\)-norm contributes
\(s^{3/2}\).  Young's inequality gives
\[
   \|e^{s\Delta}\mathbb P\operatorname{div}F\|_{L^\infty}
   \le \|K_s\|_{L^1}\|F\|_{L^\infty}
   \le Cs^{-1/2}\|F\|_{L^\infty}.
\]
The dilation \(y\mapsto e^{-a/2}y\) has \(L^\infty\)-operator norm one, and
the amplitude factor in \(S(a)\) gives \(e^{-a/2}\).  Substituting
\(s=1-e^{-a}\) proves \eqref{eq:nonlinear-semigroup-est} for smooth compactly
supported \(F\).  Since \(K_s\in L^1(\mathbb R^3)\), the convolution formula
defines \(e^{s\Delta}\mathbb P\operatorname{div}F\) pointwise almost
everywhere for every
\(F\in L^\infty(\mathbb R^3;\mathbb R^{3\times3})\), and Young's inequality
applies directly to that convolution.  Therefore the same estimate holds for
all bounded matrix fields \(F\).

Finally,
\[
   1-e^{-a}\sim a \quad (a\downarrow0),
   \qquad
   1-e^{-a}\to1 \quad (a\to\infty).
\]
Thus the integrand in \eqref{eq:semigroup-integral} is comparable to
\(a^{-1/2}\) near \(0\) and to \(e^{-a/2}\) for large \(a\), and the integral
is finite.
\end{proof}

\begin{proof}[Proof of \cref{thm:small}]
Let
\[
   M=\|V\|_{L^\infty(\mathbb R^3\times\mathbb R)}.
\]
Apply the mild formula \eqref{eq:mild} on \([\tau-T,\tau]\):
\[
   V(\tau)
   =
   S(T)V(\tau-T)
   -
   \int_0^T
   S(a)\mathbb P\operatorname{div}(V\otimes V)(\tau-a)\,da .
\]
The identity is an equality in weak-* \(L^\infty\), and the right-hand side is
represented by an \(L^\infty\) function because of
\cref{lem:stokes-estimates}.  Therefore the difference between \(V(\tau)\)
and the displayed right-hand side pairs to zero against every
\(L^1(\mathbb R^3)\) test, hence the two bounded representatives are equal
almost everywhere.  It is consequently legitimate to estimate the
\(L^\infty\)-norm of \(V(\tau)\) by the \(L^\infty\)-norm of that
representative.  The first term is bounded by \(e^{-T/2}M\).  For the Duhamel term,
\cref{lem:stokes-estimates} and \(\|V\otimes V\|_{L^\infty}\le M^2\) give
\[
   \|V(\tau)\|_{L^\infty}
   \le
   e^{-T/2}M
   +
   C_0M^2
   \int_0^T e^{-a/2}(1-e^{-a})^{-1/2}\,da .
\]
Taking the supremum in \(\tau\) gives
\[
   M\le e^{-T/2}M+C_*M^2.
\]
The solution is ancient, so \(T\) is arbitrary.  Letting \(T\to\infty\) gives
\[
   M\le C_*M^2.
\]
Choose \(\varepsilon_0=(2C_*)^{-1}\).  If \(0<M\le\varepsilon_0\), then
\(1\le C_*M\le1/2\), a contradiction.  Hence \(M=0\), and \(V\equiv0\).
\end{proof}

\section{Gaussian identity under rapid decay}

The following identity is not needed for \cref{thm:small}, but it is the
weighted energy identity naturally associated with the Gaussian measure in
self-similar variables.

\begin{proposition}[Gaussian energy identity]\label{prop:gaussian}
Let \(V\) be a smooth bounded mild ancient solution of \eqref{eq:centered},
with pressure \(P\).  Assume that, for every compact interval
\(I\subset\mathbb R\) and every integer \(m\ge0\),
\[
   \sup_{\tau\in I}\sup_{y\in\mathbb R^3}
   (1+|y|)^m\left(
   |\nabla^m V(y,\tau)|+|\nabla^m P(y,\tau)|
   \right)<\infty .
\]
Set
\[
   G(y)=e^{-|y|^2/4}.
\]
Then
\begin{equation}\label{eq:gaussian-identity}
   \frac{d}{d\tau}\int_{\mathbb R^3}|V(y,\tau)|^2G(y)\,dy
   =
   -2\int_{\mathbb R^3}|\nabla V|^2G\,dy
   -\int_{\mathbb R^3}|V|^2G\,dy
   -\frac12\int_{\mathbb R^3}(|V|^2+2P)(V\cdot y)G\,dy .
\end{equation}
\end{proposition}

\begin{proof}
All integrations by parts below may first be performed on \(B_R\) with a
cutoff \(\chi_R\in C_c^\infty(B_{2R})\), \(\chi_R=1\) on \(B_R\), \(\chi_R=0\) outside
\(B_{2R}\), and \(|\nabla^k\chi_R|\le C_kR^{-k}\).  The assumed polynomial
decay of all derivatives, together with the Gaussian factor, makes every
boundary and cutoff-error term tend to zero as \(R\to\infty\).  Thus the
formal whole-space integrations by parts are justified.

Multiply \eqref{eq:centered} by \(2VG\) and integrate over \(\mathbb R^3\).
Smoothness in \(\tau\) and dominated convergence give
\[
   2\int_{\mathbb R^3}\partial_\tau V\cdot VG\,dy
   =
   \frac{d}{d\tau}\int_{\mathbb R^3}|V|^2G\,dy .
\]
Since \(\operatorname{div} V=0\),
\[
   2\int (V\cdot\nabla)V\cdot VG
   =
   \int V\cdot\nabla(|V|^2)G
   =
   -\int |V|^2 V\cdot\nabla G.
\]
The pressure term gives
\[
   2\int \nabla P\cdot VG
   =
   -2\int P V\cdot\nabla G.
\]
For the diffusion term,
\[
\begin{aligned}
   2\int \Delta V\cdot VG
   &=
   -2\int|\nabla V|^2G
   -2\int \partial_kV_iV_i\partial_kG  \\
   &=
   -2\int|\nabla V|^2G
   -\int \partial_k(|V|^2)\partial_kG  \\
   &=
   -2\int|\nabla V|^2G
   +\int |V|^2\Delta G.
\end{aligned}
\]
Finally,
\[
   -\int (V+y\cdot\nabla V)\cdot VG
   =
   -\int |V|^2G-\frac12\int y\cdot\nabla(|V|^2)G
\]
\[
   =
   -\int |V|^2G
   +\frac12\int |V|^2(3G+y\cdot\nabla G).
\]
Using
\[
   \nabla G=-\frac y2G,\qquad
   \Delta G=\left(\frac{|y|^2}{4}-\frac32\right)G,
\]
the pure quadratic terms are
\[
   \int |V|^2\Delta G
   -\int |V|^2G
   +\frac12\int |V|^2(3G+y\cdot\nabla G)
   =
   -\int |V|^2G.
\]
The convection and pressure terms are
\[
   -\int |V|^2 V\cdot\nabla G
   -2\int P V\cdot\nabla G
   =
   \frac12\int (|V|^2+2P)(V\cdot y)G .
\]
Moving the non-time-derivative terms to the right-hand side of the equation
gives \eqref{eq:gaussian-identity}.
\end{proof}

\section{Stationary self-similar rigidity}

We now recall the stationary theorem of
Ne\v{c}as, R\r{u}\v{z}i\v{c}ka, and \v{S}verak
\cite{NRS1996}.

\begin{theorem}[Stationary self-similar rigidity]\label{thm:NRS}
Let \(W\in L^3(\mathbb R^3)\) be divergence free.  Assume that there exists
\(P\in L^{3/2}_{\mathrm{loc}}(\mathbb R^3)\) such that
\begin{equation}\label{eq:stationary-profile}
   (W\cdot\nabla)W+\nabla P
   =
   \Delta W-\frac12(W+y\cdot\nabla W)
\end{equation}
in \(\mathcal D'(\mathbb R^3)\).  Then
\[
   W\equiv0.
\]
\end{theorem}

\begin{remark}
The hypothesis \(W\in L^3(\mathbb R^3)\) is essential for this direct invocation.
Without a global integrability condition, \cref{thm:NRS} does not apply
directly.
\end{remark}

\begin{proof}[Proof of \cref{thm:stationary-limits}]
First note that each time translate \(V_n(y,\tau)=V(y,\tau+\tau_n)\) is again
a bounded mild ancient solution.  The \(L^\infty\) bound is unchanged, and
\(\operatorname{div}V_n=0\) in distributions because translation in \(\tau\)
does not affect the spatial divergence.  If \(\sigma<\tau\), applying
\eqref{eq:mild} to \(V\) on the interval
\([\sigma+\tau_n,\tau+\tau_n]\) gives
\[
   V_n(\tau)
   =
   S(\tau-\sigma)V_n(\sigma)
   -
   \int_\sigma^\tau
      S(\tau-s)\mathbb P\operatorname{div}
      (V_n\otimes V_n)(s)\,ds
\]
after the change of variables \(s\mapsto s+\tau_n\) in the weak-* time
integral.  Thus \(V_n\) satisfies the same mild formulation as \(V\).

The assumptions \eqref{eq:L3-bound-tight} and \eqref{eq:tightness} are
invariant under translations in \(\tau\).  Hence the sequence \(V_n\) has the
same \(L^\infty\), uniform \(L^3\), and uniform tightness bounds as \(V\).  By
\cref{lem:tight-closure}, the local smooth limit \(W\) satisfies the same
global \(L^3\) bound.  Since \(W\) is
independent of \(\tau\), write \(W(y,\tau)=W_0(y)\).  Then
\[
   W_0\in L^3(\mathbb R^3).
\]
We verify that this stationary function satisfies the
stationary self-similar equation in the distributional sense required by
\cref{thm:NRS}.  Since \(W\) is a
locally smooth limit of the bounded mild solutions \(V_n\), the limit-passage
argument in \cref{lem:smooth-compactness} applies to this already convergent
sequence and shows that \(W\) is itself a bounded mild ancient solution.
By \cref{lem:mild-to-projected}, \(W\) satisfies the projected
distributional equation
\[
   \partial_\tau W
   =
   \Delta W-\frac12y\cdot\nabla W-\frac12W
   -\mathbb P\operatorname{div}(W\otimes W).
\]
Because \(W(y,\tau)=W_0(y)\), its time derivative vanishes in
\(\mathcal D'(\mathbb R^3\times\mathbb R)\): for
\(\Phi\in C_c^\infty(\mathbb R^3\times\mathbb R;\mathbb R^3)\),
\[
   \langle \partial_\tau W,\Phi\rangle
   =
   -\int_{\mathbb R}\int_{\mathbb R^3}W_0(y)\cdot\partial_\tau\Phi(y,\tau)\,dy\,d\tau
   =
   -\int_{\mathbb R^3}W_0(y)\cdot
      \left(\int_{\mathbb R}\partial_\tau\Phi(y,\tau)\,d\tau\right)dy
   =
   0 .
\]
For the stationary velocity \(W(y,\tau)=W_0(y)\), the Leray pressure
\[
   P_0=R_iR_j(W_{0,i}W_{0,j})
\]
belongs to \(L^{3/2}(\mathbb R^3)\), because \(W_0\in L^3\), hence
\(W_{0,i}W_{0,j}\in L^{3/2}(\mathbb R^3)\), and the Riesz transforms are
bounded on \(L^{3/2}\) \cite{Stein1970}.  In
\(\mathcal D'(\mathbb R^3)\),
\[
   (I-\mathbb P)\operatorname{div}(W_0\otimes W_0)=-\nabla P_0,
   \qquad
   \mathbb P\nabla P_0=0.
\]
This is the Helmholtz decomposition in the whole space
\cite{Sohr2001,Galdi2011}.  Hence
\[
   \mathbb P\operatorname{div}(W_0\otimes W_0)
   =
   \operatorname{div}(W_0\otimes W_0)+\nabla P_0 .
\]
Substituting this identity into the projected stationary equation converts it
into
\[
   \operatorname{div}(W_0\otimes W_0)+\nabla P_0
   =
   \Delta W_0-\frac12(W_0+y\cdot\nabla W_0)
\]
in \(\mathcal D'(\mathbb R^3)\).

For every test function \(\varphi\in C_c^\infty(\mathbb R^3;\mathbb R^3)\),
testing the resulting distributional identity against
\(\psi(\tau)\varphi(y)\), where \(\psi\in C_c^\infty(\mathbb R)\) and
\(\int\psi\,d\tau=1\), leaves the same spatial weak form because all terms in
the identity are independent of \(\tau\).  Finally,
\(\operatorname{div}W_0=0\) gives, pointwise because \(W_0\) is smooth,
\[
   \operatorname{div}(W_0\otimes W_0)_i
   =
   \partial_j(W_{0,i}W_{0,j})
   =
   W_{0,j}\partial_jW_{0,i}+W_{0,i}\partial_jW_{0,j}
   =
   (W_0\cdot\nabla)W_{0,i}
\]
and hence the same identity holds in \(\mathcal D'(\mathbb R^3)\).  Therefore
\[
   (W_0\cdot\nabla)W_0+\nabla P_0
   =
   \Delta W_0-\frac12(W_0+y\cdot\nabla W_0),
   \qquad \operatorname{div}W_0=0 .
\]
Thus \(W_0\) is precisely a stationary self-similar profile in the class
covered by \cref{thm:NRS}.  Consequently \(W_0\equiv0\), and hence
\(W\equiv0\).
\end{proof}

\begin{proof}[Proof of \cref{cor:stationary-tight}]
Apply \cref{thm:stationary-limits} with \(\tau_n\equiv0\).
\end{proof}

\begin{proof}[Proof of \cref{cor:size-threshold}]
This is the contrapositive of \cref{thm:small}.
\end{proof}

\section{Concluding remarks}

These results do not by themselves imply a Type I exclusion theorem.
They prove two conditional Liouville statements for bounded mild ancient
solutions of \eqref{eq:centered}.  A stationary local limit is zero whenever
uniform \(L^3\)-tightness supplies the global \(L^3\) hypothesis in the
Ne\v{c}as--R\r{u}\v{z}i\v{c}ka--\v{S}verak theorem.  Independently, a mild ancient solution with
sufficiently small \(L^\infty(\mathbb R^3\times\mathbb R)\) norm is zero.  Thus any
nontrivial bounded mild ancient solution must carry a definite
\(L^\infty\)-size, and any argument which produces a stationary limit must
also provide enough spatial compactness to retain the critical \(L^3\) bound.

\begin{thebibliography}{99}

\bibitem{ESS2003}
L. Escauriaza, G. A. Seregin, and V. \v{S}verak,
\emph{\(L_{3,\infty}\)-solutions of Navier--Stokes equations and backward uniqueness},
Russian Math. Surveys \textbf{58} (2003), no. 2, 211--250,
doi:10.1070/RM2003v058n02ABEH000609.

\bibitem{Galdi2011}
G. P. Galdi,
\emph{An introduction to the mathematical theory of the Navier--Stokes
equations},
2nd ed., Springer Monographs in Mathematics, Springer, New York, 2011.

\bibitem{GilbargTrudinger2001}
D. Gilbarg and N. S. Trudinger,
\emph{Elliptic partial differential equations of second order},
Classics in Mathematics, Springer, Berlin, 2001.

\bibitem{GigaMiyakawa1985}
Y. Giga and T. Miyakawa,
\emph{Solutions in \(L_r\) of the Navier--Stokes initial value problem},
Arch. Rational Mech. Anal. \textbf{89} (1985), no. 3, 267--281,
doi:10.1007/BF00276875.

\bibitem{Kato1984}
T. Kato,
\emph{Strong \(L^p\)-solutions of the Navier--Stokes equation in
\(\mathbb R^m\), with applications to weak solutions},
Math. Z. \textbf{187} (1984), no. 4, 471--480.

\bibitem{Ladyzhenskaya1968}
O. A. Ladyzhenskaya, V. A. Solonnikov, and N. N. Ural'tseva,
\emph{Linear and quasilinear equations of parabolic type},
Translations of Mathematical Monographs, vol. 23, American Mathematical
Society, Providence, RI, 1968.

\bibitem{NRS1996}
J. Ne\v{c}as, M. R\r{u}\v{z}i\v{c}ka, and V. \v{S}verak,
\emph{On Leray's self-similar solutions of the Navier--Stokes equations},
Acta Math. \textbf{176} (1996), no. 2, 283--294,
doi:10.1007/BF02551584.

\bibitem{Pazy1983}
A. Pazy,
\emph{Semigroups of linear operators and applications to partial differential
equations},
Applied Mathematical Sciences, vol. 44, Springer, New York, 1983.

\bibitem{Sohr2001}
H. Sohr,
\emph{The Navier--Stokes equations: an elementary functional analytic
approach},
Birkh\"auser Advanced Texts, Birkh\"auser, Basel, 2001.

\bibitem{Stein1970}
E. M. Stein,
\emph{Singular integrals and differentiability properties of functions},
Princeton University Press, 1970.

\end{thebibliography}

\end{document}
